% =====================================================================
%  Compte rendu de projet
%  Module : Développement et déploiement d'applications intelligentes
%  Modèle à suivre obligatoirement. Compiler avec : latexmk -pdf main.tex
% =====================================================================
\documentclass[11pt,a4paper]{article}

% ---------- Encodage, langue, police ----------
\usepackage[utf8]{inputenc}
\usepackage[T1]{fontenc}
\usepackage{lmodern}
\usepackage[french]{babel}
\usepackage{csquotes}

% ---------- Mise en page ----------
\usepackage[margin=2.5cm]{geometry}
\usepackage{fancyhdr}
\usepackage{enumitem}
\setlength{\parskip}{0.5em}

% ---------- Tableaux, figures, code ----------
\usepackage{graphicx}
\usepackage{booktabs}
\usepackage{tabularx}
\usepackage{array}
\usepackage{caption}
\usepackage{float}
\usepackage[table]{xcolor}
\usepackage{listings}
\usepackage[most]{tcolorbox}
\usepackage[round]{natbib}
\usepackage[hidelinks]{hyperref}

% =====================================================================
%  1. INFORMATIONS DU PROJET : à remplir par l'équipe
% =====================================================================
\newcommand{\titreprojet}{Titre du projet}
\newcommand{\numeroprojet}{N° du catalogue (ou hors catalogue)}
\newcommand{\nomequipe}{Équipe XX}
\newcommand{\depot}{https://github.com/organisation/depot-equipe}
\newcommand{\versionfinale}{v1.0 (étiquette Git de la version rendue)}
\newcommand{\enseignant}{Nom de l'enseignant}
\newcommand{\daterendu}{JJ/MM/AAAA}

% Une ligne par membre (2 à 4) : Nom \& Prénom & compte GitHub & rôle principal \\
\newcommand{\membres}{%
  Nom Prénom 1 & @github1 & Rôle principal \\
  Nom Prénom 2 & @github2 & Rôle principal \\
  % Nom Prénom 3 & @github3 & Rôle principal \\
  % Nom Prénom 4 & @github4 & Rôle principal \\
}

% =====================================================================
%  2. CONSIGNES : passer à \consignesfalse pour la version rendue
% =====================================================================
\newif\ifconsignes
\consignestrue
% \consignesfalse

\newtcolorbox{boiteconsigne}{colback=blue!4, colframe=blue!50!black,
  fonttitle=\bfseries\small, fontupper=\small, title=Consigne, breakable}
\newcommand{\consigne}[1]{\ifconsignes\begin{boiteconsigne}#1\end{boiteconsigne}\fi}

% Zone à compléter, en rouge pour qu'aucune ne soit oubliée
\newcommand{\acompleter}[1][à compléter]{\textcolor{red}{\textbf{[#1]}}}

% Rédacteur(s) de la section : sert de preuve de contribution
\newcommand{\redacteur}[1]{\noindent\textit{\small Rédaction : #1}\par\medskip}

% Figure ou emplacement réservé si le fichier n'existe pas encore
% Usage : \figure{fichier}{largeur}{légende}{étiquette}
\newcommand{\insererfigure}[4]{%
  \begin{figure}[H]\centering
  \IfFileExists{figures/#1}%
    {\includegraphics[width=#2\linewidth]{figures/#1}}%
    {\fbox{\parbox[c][4cm][c]{#2\linewidth}{\centering\textcolor{red}{Figure manquante : figures/#1}}}}%
  \caption{#3}\label{#4}\end{figure}}

% ---------- Style du code ----------
\lstdefinestyle{code}{
  basicstyle=\ttfamily\footnotesize, numbers=left, numberstyle=\tiny\color{gray},
  keywordstyle=\color{blue!70!black}\bfseries, commentstyle=\color{green!45!black},
  stringstyle=\color{orange!70!black}, frame=single, rulecolor=\color{gray!40},
  backgroundcolor=\color{gray!4}, breaklines=true, showstringspaces=false,
  captionpos=b, tabsize=4,
  literate={é}{{\'e}}1 {è}{{\`e}}1 {ê}{{\^e}}1 {à}{{\`a}}1 {ç}{{\c c}}1
           {ù}{{\`u}}1 {ô}{{\^o}}1 {î}{{\^i}}1 {É}{{\'E}}1 {«}{{\guillemotleft}}1 {»}{{\guillemotright}}1
}
\lstset{style=code}
\renewcommand{\lstlistingname}{Extrait}

% ---------- En-tête et pied de page ----------
\pagestyle{fancy}
\fancyhf{}
\fancyhead[L]{\small\nomequipe}
\fancyhead[R]{\small\titreprojet}
\fancyfoot[C]{\thepage}

% =====================================================================
\begin{document}

% Page de garde : ne pas modifier la mise en page, seulement les informations dans main.tex
\begin{titlepage}
\centering
{\large Développement et déploiement d'applications intelligentes\par}
\vspace{0.5cm}
{\normalsize Compte rendu de projet\par}
\vspace{2.5cm}
{\LARGE\bfseries \titreprojet\par}
\vspace{0.5cm}
{\large Projet \numeroprojet\par}
\vspace{2cm}
{\large\bfseries \nomequipe\par}
\vspace{0.5cm}
\begin{tabular}{lll}
\toprule
\textbf{Membre} & \textbf{GitHub} & \textbf{Rôle principal} \\
\midrule
\membres
\bottomrule
\end{tabular}
\vfill
\begin{tabular}{rl}
Dépôt : & \url{\depot} \\
Version rendue : & \versionfinale \\
Enseignant : & \enseignant \\
Date : & \daterendu \\
\end{tabular}
\end{titlepage}

\section*{Résumé}
\consigne{10 lignes au plus. Le problème, la solution construite, les technologies principales et le résultat chiffré le plus important. Le résumé doit pouvoir être lu seul. Rédigez-le en dernier.}

\acompleter

\medskip
\noindent\textbf{Mots-clés :} \acompleter[4 à 6 mots-clés]
\newpage


\tableofcontents
\newpage

\section{Introduction}
\redacteur{\acompleter[nom]}

\consigne{1 page environ. Reprenez et améliorez le livrable 1 du dossier de lancement (sections 1.1 à 1.3). Terminez par l'annonce du plan.}

\subsection{Contexte et problème}
\acompleter

\subsection{Objectif}
\acompleter[objectif en une phrase]

\subsection{Utilisateurs cibles et scénario d'usage}
\acompleter

\subsection{Plan du rapport}
\acompleter

\section{Cahier des charges}
\redacteur{\acompleter[nom]}

\consigne{Présentez le MVP tel que validé en semaine 1, puis ce qui a réellement été livré. Tout écart doit être expliqué : un écart justifié n'est pas pénalisé, un écart caché l'est (critère P1 de la grille).}

\subsection{Fonctionnalités du MVP}

\begin{table}[H]\centering
\caption{Fonctionnalités prévues et état à la livraison}\label{tab:mvp}
\begin{tabularx}{\linewidth}{lXll}
\toprule
\textbf{N°} & \textbf{Fonctionnalité} & \textbf{Prévue en S1} & \textbf{État final} \\
\midrule
F1 & \acompleter & oui & livrée et testée \\
F2 & \acompleter & oui & \acompleter \\
F3 & \acompleter & oui & \acompleter \\
\bottomrule
\end{tabularx}
\end{table}

\subsection{Exigence liée à la taille de l'équipe}
\consigne{Trinôme : composant supplémentaire (évaluation automatisée, monitoring ou second agent). Équipe de 4 : deux agents ou deux sous-systèmes intégrés. Binôme : écrire « sans objet ».}
\acompleter

\subsection{Hors périmètre et écarts avec la semaine 1}
\acompleter

\section{Données}
\redacteur{\acompleter[nom]}

\consigne{Reprenez le livrable 2 et complétez-le avec ce que vous avez réellement fait : prétraitement, nettoyage, découpage. Citez la source et sa licence dans la bibliographie.}

\subsection{Source, licence et description}

\begin{table}[H]\centering
\caption{Description du jeu de données}\label{tab:donnees}
\begin{tabularx}{\linewidth}{lX}
\toprule
Source & \acompleter{} \citep{jeu-de-donnees} \\
Licence & \acompleter \\
Volume & \acompleter \\
Format & \acompleter \\
Données personnelles & \acompleter[aucune, ou traitement appliqué] \\
\bottomrule
\end{tabularx}
\end{table}

\subsection{Prétraitement}
\acompleter[étapes, outils, choix (taille des morceaux, filtrage, normalisation)]

\subsection{Jeu de test}
\consigne{Taille, origine, mode de construction des réponses attendues. Le jeu de test complet va en annexe ou dans le dépôt.}
\acompleter

\section{Architecture}
\redacteur{\acompleter[nom]}

\consigne{Le schéma est obligatoire. Partez du schéma Mermaid du livrable 3, exportez-le en PDF ou PNG (par exemple sur mermaid.live) et placez-le dans \texttt{figures/}. Chaque composant du schéma est expliqué dans le texte et chaque choix technique est justifié face à au moins une alternative.}

\insererfigure{architecture.pdf}{0.95}{Architecture du système}{fig:architecture}

\subsection{Composants}

\begin{table}[H]\centering
\caption{Composants et choix techniques}\label{tab:composants}
\begin{tabularx}{\linewidth}{lllX}
\toprule
\textbf{Composant} & \textbf{Technologie} & \textbf{Alternative écartée} & \textbf{Justification} \\
\midrule
Modèle & \acompleter & \acompleter & \acompleter \\
Recherche ou mémoire & \acompleter{} & \acompleter{} & \acompleter{} (voir par exemple \citet{lewis2020rag}) \\
API & FastAPI \citep{fastapi} & \acompleter & \acompleter \\
Interface & \acompleter & \acompleter & \acompleter \\
\bottomrule
\end{tabularx}
\end{table}

\subsection{Parcours d'une requête}
\acompleter[de l'action de l'utilisateur jusqu'à la réponse, en suivant la figure~\ref{fig:architecture}]

\subsection{Agents et orchestration}
\consigne{Rôle de chaque agent, outils qu'il peut appeler, façon dont les agents échangent. Obligatoire pour une équipe de 4.}
\acompleter

\section{Implémentation}
\redacteur{\acompleter[nom]}

\consigne{Ne recopiez pas votre code : le dépôt est là pour ça. Montrez l'organisation du code, les points d'entrée de l'API et deux ou trois extraits courts (moins de 20 lignes) qui éclairent un choix important. Les prompts complets vont en annexe.}

\subsection{Organisation du dépôt}
\acompleter[arborescence commentée des dossiers principaux]

\subsection{API}

\begin{table}[H]\centering
\caption{Points d'entrée de l'API}\label{tab:api}
\begin{tabularx}{\linewidth}{llX}
\toprule
\textbf{Méthode} & \textbf{Route} & \textbf{Rôle} \\
\midrule
POST & \texttt{/ask} & \acompleter \\
GET & \texttt{/health} & Vérification de l'état du service \\
\bottomrule
\end{tabularx}
\end{table}

\subsection{Extraits significatifs}

\begin{lstlisting}[language=Python, caption={Exemple : route principale de l'API (à remplacer)}, label={lst:api}]
@app.post("/ask")
def ask(question: Question) -> Answer:
    passages = retriever.search(question.text, k=5)
    return agent.answer(question.text, passages)
\end{lstlisting}

\acompleter[explication de l'extrait~\ref{lst:api}]

\section{Qualité logicielle et déploiement}
\redacteur{\acompleter[nom]}

\consigne{Correspond aux critères P4 et P5 de la grille. Donnez des faits vérifiables : nombre de tests, couverture mesurée, lien vers la dernière exécution de la CI, commande de lancement.}

\subsection{Tests}
\acompleter[types de tests, outils, nombre, couverture]

\subsection{Intégration continue}
\acompleter[étapes du workflow GitHub Actions, lien vers une exécution réussie]

\subsection{Conteneurisation et lancement}

\begin{lstlisting}[language=bash, caption={Lancement de l'application}]
docker build -t nom-projet .
docker run -p 8000:8000 --env-file .env nom-projet
\end{lstlisting}

\subsection{Suivi des expériences et des données}
\consigne{Si vous utilisez MLflow ou DVC, montrez ce qui est suivi et comment. Sinon, écrivez « non utilisé » et pourquoi.}
\acompleter

\section{Évaluation}
\redacteur{\acompleter[nom]}

\consigne{Section la plus importante pour le critère P3. Mesurez chaque critère de réussite défini en semaine 1 sur le jeu de test, comparez à la cible, puis analysez les erreurs avec des exemples concrets. Un résultat inférieur à la cible, bien analysé, vaut mieux qu'un bon chiffre non expliqué.}

\subsection{Résultats}

\begin{table}[H]\centering
\caption{Critères de réussite : cibles et résultats}\label{tab:resultats}
\begin{tabularx}{\linewidth}{Xlll}
\toprule
\textbf{Indicateur} & \textbf{Cible} & \textbf{Résultat} & \textbf{Atteint} \\
\midrule
\acompleter & \acompleter & \acompleter & oui / non \\
\acompleter & \acompleter & \acompleter & oui / non \\
\acompleter & \acompleter & \acompleter & oui / non \\
\bottomrule
\end{tabularx}
\end{table}

\subsection{Protocole}
\acompleter[comment les mesures ont été faites : jeu de test, nombre d'exécutions, matériel, modèle et paramètres]

\subsection{Analyse des erreurs}
\consigne{Classez les échecs par type (mauvaise recherche, hallucination, mauvais outil appelé, délai\ldots), donnez leur fréquence et au moins deux exemples commentés.}
\acompleter

\section{Sécurité, risques et éthique}
\redacteur{\acompleter[nom]}

\consigne{Reprenez les risques déclarés en semaine 1 et dites ce que sont devenues les parades. Critère P6 de la grille.}

\begin{table}[H]\centering
\caption{Risques et parades}\label{tab:risques}
\begin{tabularx}{\linewidth}{XXl}
\toprule
\textbf{Risque} & \textbf{Parade mise en place} & \textbf{Testée ?} \\
\midrule
\acompleter & \acompleter & oui / non \\
\acompleter & \acompleter & oui / non \\
\bottomrule
\end{tabularx}
\end{table}

\subsection{Garde-fous sur les actions de l'agent}
\acompleter[liste blanche, validation humaine, journalisation, ou « sans objet »]

\subsection{Gestion des secrets et des données personnelles}
\acompleter

\subsection{Limites et usages à éviter}
\acompleter

\section{Gestion de projet}
\redacteur{\acompleter[nom]}

\consigne{Critère P7 et preuve pour le coefficient individuel. Le tableau des contributions doit être cohérent avec les rapports hebdomadaires et l'historique Git. Chaque membre relit et valide sa propre ligne.}

\subsection{Planning prévu et réel}

\begin{table}[H]\centering
\caption{Jalons prévus et atteints}\label{tab:planning}
\begin{tabularx}{\linewidth}{lXXl}
\toprule
\textbf{Semaine} & \textbf{Jalon prévu} & \textbf{Réalisé} & \textbf{Écart} \\
\midrule
1 & Dossier de lancement, prototype en script & \acompleter & \acompleter \\
2 & Agent exposé par l'API, premiers tests & \acompleter & \acompleter \\
3 & MVP complet, Docker et CI & \acompleter & \acompleter \\
4 & Évaluation, gel du code, soutenance & \acompleter & \acompleter \\
\bottomrule
\end{tabularx}
\end{table}

\subsection{Contributions}

\begin{table}[H]\centering
\caption{Contribution de chaque membre}\label{tab:contributions}
\begin{tabularx}{\linewidth}{lXXl}
\toprule
\textbf{Membre} & \textbf{Fonctionnalités et parties réalisées} & \textbf{Sections du rapport} & \textbf{Part (\%)} \\
\midrule
\acompleter & \acompleter & \acompleter & \acompleter \\
\acompleter & \acompleter & \acompleter & \acompleter \\
\bottomrule
\end{tabularx}
\end{table}

\subsection{Difficultés rencontrées}
\acompleter[principales difficultés et comment elles ont été résolues]

\section{Conclusion et perspectives}
\redacteur{\acompleter[nom]}

\consigne{Une demi-page à une page. Bilan honnête par rapport à l'objectif, principales limites, deux ou trois pistes d'amélioration concrètes et classées par priorité. Pas de nouveau résultat ici.}

\subsection{Bilan}
\acompleter

\subsection{Limites}
\acompleter

\subsection{Perspectives}
\acompleter


\bibliographystyle{plainnat}
\bibliography{references}

\appendix
\section{Installation et reproduction}
\consigne{Les commandes exactes pour cloner, configurer et lancer le projet sur une machine neuve. L'enseignant les exécutera.}

\begin{lstlisting}[language=bash]
git clone URL-DU-DEPOT
cd depot
cp .env.example .env   # renseigner les cles
docker compose up --build
\end{lstlisting}

\section{Prompts}
\consigne{Les prompts système et les modèles de prompts utilisés, en version finale.}
\acompleter

\section{Extrait du jeu de test}
\acompleter[10 cas représentatifs avec réponse attendue et réponse obtenue]

\section{Déclaration d'utilisation d'outils d'IA générative}
\consigne{Obligatoire. Indiquez quels outils d'IA (assistants de code, LLM conversationnels) ont servi, pour quelles parties du code ou du rapport, et comment les résultats ont été vérifiés. Écrivez « aucun » le cas échéant.}

\begin{table}[H]\centering
\caption{Outils d'IA générative utilisés}\label{tab:ia}
\begin{tabularx}{\linewidth}{lXX}
\toprule
\textbf{Outil} & \textbf{Usage} & \textbf{Vérification effectuée} \\
\midrule
\acompleter & \acompleter & \acompleter \\
\bottomrule
\end{tabularx}
\end{table}


\end{document}
